\documentclass[titlepage]{ctexart}

\usepackage{amsmath}
\usepackage{amssymb}
\usepackage{caption}

\title{第一章~数列}
\author{dovego}
\date{\today}

\begin{document}
\maketitle


\section{数列的概念及其函数特性}

\subsection{数列的概念}

按一定次序排列的一列数叫做\textbf{数列}, 数列中的每一个数叫做这个数列的\textbf{项}, 数列的一般形式可以表示为
\[
a_{1}, a_{2}, a_{3}, \cdots, a_{n}, \cdots
\]
也可以简写为 $\{a_{n}\}$. 其中 $a_{1}$ 是数列的第一项, 也叫\textbf{首项}; $a_{n}$ 是数列的第 $n$ 项, 也叫\textbf{通项}.

项数有限的数列叫做\textbf{有穷数列}, 项数无限的数列叫做\textbf{无穷数列}.

数列也可以看作定义域在正整数集 $\mathbf{N_+}$ (或其子集)上的函数.

若数列 $\{a_{n}\}$ 的第 $n$ 项 $a_{n}$ 与 $n$ 之间的函数关系可以用一个式子 $a_{n} = f(n)$ 来表示, 那么这个式子就叫做这个数列的\textbf{通项公式}.
注: (1) 不是所有数列都能写出通项公式; (2) 通项公式的形式不唯一.

\subsection{数列的函数特性}

\textbf{递增数列}: 对所有 $n \geqslant 1(n \in \mathbf{N_+})$, 都有 $a_{n + 1} > a_{n}$.

\textbf{递减数列}: 对所有 $n \geqslant 1$, 都有 $a_{n+1} > a_{n}$.

\textbf{常数列}: 各项都相等.


\section{等差数列}

\subsection{等差数列的概念及其通项公式}

\textbf{等差数列}: 对所有 $n \geqslant 1(n \in \mathbf{N_+})$, 都有 $a_{n+1} - a_{n} = d(d\text{为常数})$. 这个常数 $d$ 称为\textbf{公差}.

用\textbf{累加法}可以得出等差数列 $\{a_{n}\}$ 的通项公式为: $a_{n} = a_{1} + (n-1)d, n = 1,2,3,\dots$.

从函数的角度来看, 当 $d > 0$ 时, $\{a_{n}\}$ 为递增数列; 当 $d < 0$ 时, $\{a_{n}\}$ 为递减数列; 当 $d = 0$ 时, $\{a_{n}\}$为常数列.

\textbf{等差中项}: 若在 $a, b$ 中间插入一个数 $A$, 使$a, A, b$ 构成等差数列, 那么 $A$ 叫做 $a$ 和 $b$ 的\textbf{等差中项}. 且 $A = \dfrac{a+b}{2}$.

\textbf{补充} \quad 对于等差数列 $\{a_n\}$, 有 $\dfrac{a_m + a_n}{2} = a_{\frac{m+n}{2}}$.

\subsection{等差数列的前 $n$ 项和}

由\textbf{倒序相加法}可以得出等差数列的前 $n$ 项和公式为
\[
S_n = \frac{n(a_1 + a_n)}{2} = na_1 + \frac{n(n-1)}{2}d.
\]


\section{等比数列}

\subsection{等比数列的概念及其通项公式}

\textbf{等比数列}: 对所有 $n \geqslant 1(n \in \mathbf{N_+})$, 都有 $\dfrac{a_{n+1}}{a_n} = q(q\text{为常数且} q \neq 0)$. 这个常数 $q$ 称为\textbf{公比}.

用\textbf{累乘法}可以得出等比数列的通项公式为: $a_n = a_1 q^{n-1} (a_1 \neq 0, q \neq 0)$.

\textbf{等比中项}: 若在 $a, b$ 中插入一个数 $G$, 使 $a, G, b$ 构成等比数列, 那么 $G$ 叫做 $a$ 和 $b$ 的\textbf{等比中项}. 且 $G^2 = ab, G = \pm \sqrt{ab}$. 

\textbf{补充} \quad 对于等比数列 $\{a_n\}$, 有 $a_ma_n = a_{\frac{m+n}{2}}^2$.

\subsection{等比数列的前 $n$ 项和}

由\textbf{错位相减法}可以得出等比数列的前 $n$ 项和公式为
\[
S_n = 
\begin{cases}
    na_1, &q= 1, \\
    \dfrac{a_1(1-q^n)}{1-q}, &q\neq 1 \text{且}q \neq 0.
\end{cases}
\]


\section{数列在日常经济生活中的应用}

单利, 复利, 零存整取模型, 定期自动转存模型, 分期付款模型.


\section{数学归纳法}

\textbf{数学归纳法}是用来证明某些与正整数 $n$ 有关的数学命题的一种方法. 它的基本步骤是:

(1) 证明: 当 $n$ 取第一个值 $n_0$($n_0$ 是一个确定的正整数, 如 $n_0 = 1$ 或 2 等)时, 命题成立;

(2) 假设当 $n = k(k \in \mathbf{N_+}, k \geqslant n_0$ 时命题成立, 证明当 $n \geqslant k$ 时, 命题也成立. 

根据(1)(2)可以断定命题对一切从 $n_0$ 开始的正整数 $n$ 都成立.




\end{document}